\section{Local magnetic subspace and spin vertices}
\label{sec:local-torque}

\subsection{Hermitian local partition}

Let $P_\ell$ project onto the explicitly selected localized magnetic Wannier orbitals of site $\ell$, with $P_M=\sum_\ell P_\ell$. The orthonormal limit of the Hermitian local operator proposed in Ref.~\cite{MartinezCarracedo2023} is
\begin{equation}
\bL_\ell[X]=\frac{1}{2}(P_\ell X+XP_\ell).
\label{eq:local-partition}
\end{equation}
The local exchange fields are
\begin{equation}
H_{\XC,\ell}=\bL_\ell[H_{\XC}],
\qquad
\sum_\ell H_{\XC,\ell}
=\frac12(P_MH_{\XC}+H_{\XC}P_M).
\label{eq:local-xc}
\end{equation}
The difference between $H_{\XC}$ and the right-hand side is an explicit magnetic-support residual. The onsite-only choice $P_\ell H_{\XC}P_\ell$ is useful as a sensitivity test but is not the default.

The definition of $P_\ell$ is physical rather than merely numerical. Ref.~\cite{MartinezCarracedo2023} reports large changes when nonlocalized orbitals are included in the rotated subspace. A production calculation must therefore list the selected Wannier functions and compare at least a localized magnetic-orbital set with a larger atom-centered set.

\subsection{Local frames and transverse coordinates}

At each magnetic site, choose a right-handed frame
\begin{equation}
(\bm t_{\ell1},\bm t_{\ell2},\bm n_\ell),
\qquad
\bm t_{\ell1}\times\bm t_{\ell2}=\bm n_\ell,
\label{eq:local-frame}
\end{equation}
where $\bm n_\ell$ is the equilibrium moment direction. A small unit-vector fluctuation is
\begin{equation}
\bm e_\ell=\bm n_\ell
+\pi_{\ell1}\bm t_{\ell1}
+\pi_{\ell2}\bm t_{\ell2}+O(\pi^2).
\label{eq:pi-coordinate}
\end{equation}
The transverse-direction vertex is
\begin{equation}
\bT_{\ell a}
=\left.\frac{\partial H}{\partial\pi_{\ell a}}\right|_0.
\label{eq:direction-vertex}
\end{equation}
For a rigid local exchange field
\begin{equation}
H_{\XC,\ell}
=\frac12\Delta_\ell\otimes
(\bm n_\ell\cdot\bm\sigma),
\end{equation}
this becomes
\begin{equation}
\bT_{\ell a}
=\frac12\Delta_\ell\otimes
(\bm t_{\ell a}\cdot\bm\sigma).
\label{eq:rigid-direction-vertex}
\end{equation}

\subsection{Rotation-angle vertex and exchange-field-only rotation}

For an infinitesimal rotation around a unit axis $\bm u_{\ell c}$, the first-order exchange-field perturbation is~\cite{MartinezCarracedo2023}
\begin{equation}
\Gamma_{\ell c}
=\frac{\ii}{2}
[H_{\XC,\ell},\bm u_{\ell c}\cdot\bm\sigma]
=-\ii[S_{\ell c},H_{\XC,\ell}].
\label{eq:rotation-vertex}
\end{equation}
Only $H_{\XC,\ell}$ is rotated. Commuting the generator with the full Hamiltonian would rotate SOC relative to the lattice and does not implement the intended DFT-to-spin mapping~\cite{MartinezCarracedo2023}.

Direction and angle derivatives differ by a tangent-plane rotation. With $\bm u_{\ell1}=\bm t_{\ell1}$ and $\bm u_{\ell2}=\bm t_{\ell2}$,
\begin{equation}
\Gamma_{\ell1}=-\bT_{\ell2},
\qquad
\Gamma_{\ell2}=+\bT_{\ell1}.
\label{eq:coordinate-map}
\end{equation}
This distinction is essential for circular magnon combinations and torque signs.

\subsection{Finite-q representation of the local partition}

The real-space local partition generally produces a vertex that depends on both $\kk$ and $\qq$. Introduce the q-dependent site selector
\begin{equation}
P_\ell(\qq)
=\left\langle w_{\kk+\qq}\middle|
\sum_{\rr}e^{\ii\qq\cdot(\rr+\bm\tau_\ell)}P_{\rr\ell}
\middle|w_{\kk}\right\rangle.
\label{eq:q-projector}
\end{equation}
For diagonal orbital projectors in the atomic-position gauge,
\begin{equation}
[P_\ell(\qq)]_{nm}
=\delta_{nm}\delta_{s(n),\ell}
 e^{\ii\qq\cdot(\bm\tau_\ell-\bm\tau_n)}.
\label{eq:q-projector-atomic}
\end{equation}
Let $\Gamma_c^{\rm global}(\kk)$ be the uniform exchange-field rotation around the local axis selected for site $\ell$. Fourier transformation of the Hermitian local partition gives the project implementation rule
\begin{equation}
\Gamma_{\ell c}(\kk+\qq,\kk)
=\frac12\left[
P_\ell(\qq)\Gamma_c^{\rm global}(\kk)
+\Gamma_c^{\rm global}(\kk+\qq)P_\ell(\qq)
\right].
\label{eq:q-local-vertex}
\end{equation}
Only in the strictly onsite atomic-center limit is this matrix independent of $\kk$ and $\qq$. Equation~\eqref{eq:q-local-vertex} is the orthonormal finite-q implementation of the local-partition prescription and is classified as a project extension.
